% !TEX root = ../paper.tex
\subsection{Precomputed Ranks for Ordering}
\label{sec:rank}

\Cref{piop:lex-sorted-check} inspects the characters of the sorted column, one position per round. In a relational plan, however, the strings being sorted are almost never new: every operator between a base table and the \texttt{ORDER BY} is data- or domain-preserving (\cref{sec:consistency-checking}), so each string in the sorted column is a copy of a string in a committed base column, and the \emph{relative order} of those strings is a property of the base column alone. We can therefore compute that order once, at ingestion, and store it next to each string. This is the rank column $\mathsf{rk}$ of \cref{sec:string-arith}: $\mathsf{rk}[i]$ is the dense rank of the $i$-th string among the distinct strings of its column, so that
\begin{displaymath}
  \begin{aligned}
    \mathsf{rk}[i] < \mathsf{rk}[j] &\iff s_i \lexlt s_j, \\
    \mathsf{rk}[i] = \mathsf{rk}[j] &\iff s_i = s_j .
  \end{aligned}
\end{displaymath}
Like the hash and the length, the rank is a str-level column inherited through every domain-preserving update by the provenance lookup of \cref{piop:domain-preserving-update-check}, so it is available, and trusted, wherever the string itself is. Two remarks on the definition. First, the rank must be \emph{dense}: if equal strings received distinct ranks, two rows with the same string would be forced into a fixed relative order instead of being ordered by the next sort key, which breaks the multi-key semantics of \cref{sec:sql-string-order}. Second, ranks are only comparable within one base column; the plan's lineage tells the planner which base column a sort key descends from, and two keys that descend from the same base column (such as the two copies of \texttt{n\_name} that Q7 of TPC-H joins and orders by) share a rank space.

\parhead{Rank validity}
A rank column is valid when it is monotone in the lexicographic order and separates exactly the distinct strings. Since $\lexle$ is a total preorder, it suffices to check this on the column \emph{sorted by rank}: writing $\pi$ for the permutation that sorts $\mathsf{rk}$ non-decreasingly, the relation $\NPRelation_{\mathsf{RankValid}}$ (stated in \cref{sec:proof-rankvalid}) requires that along $\pi$ the strings be non-decreasing in $\lexle$ and that two adjacent ranks be equal exactly when the two adjacent strings are equal. Adjacent consistency along $\pi$ implies the global property by transitivity: ranks are non-decreasing along the lexicographic order and increase exactly where the string changes.

\parhead{Certifying the ranks}
As discussed at the start of this section, the data owner may simply commit the rank column alongside the hash column and the fingerprint bins as part of the data. If instead the ranks are to be \emph{verified}, the char-level protocol of \cref{sec:lex-sorted-check} is exactly the tool: the prover materializes the base column sorted by rank, proves it is a permutation of the base column, proves the ranks non-decreasing along it, proves it lexicographically sorted, and finally proves that rank ties coincide with string ties, using the equal-neighbour indicator returned by \cref{piop:lex-sorted-check}, which marks precisely the adjacent pairs of equal strings. This is done once per committed string column, at ingestion, and amortizes over every query that later sorts a descendant of that column.

\begin{DefinePIOP}{Rank Certification Check}{rank-cert}[Rank Certification Check]
  \emergencystretch=3em\relax
  \emph{Inputs:} the committed column $\Oracle{\MLE{\strcol}}$ with its rank oracle $\Oracle{\MLE{\mathsf{rk}}}$; and $\Oracle{\MLE{\mathsf{ind}}}$.
  \vspace{2mm}
  \begin{enumerate}[nosep]
    \item (\emph{Sorted copy})\label[step]{rc:copy} $\PIOPProver$ materializes the column $\strcol^{\pi}$ obtained by sorting $\strcol$ by rank, with its rank column $\mathsf{rk}^{\pi}$, and sends all of its oracles. $\PIOPProver$ and $\PIOPVerifier$ invoke \cref{piop:domain-preserving-update-check} on $(\strcol, \strcol^{\pi})$, and \cref{piop:mset} on $(\mathsf{a}^{\pi}, \mathsf{h}^{\pi})$ and $(\mathsf{a}, \mathsf{h})$, as in \cref{ss:valid,ss:perm}.
    \item (\emph{Ranks non-decreasing})\label[step]{rc:mono} $\PIOPProver$ sends $\Oracle{\MLE{\mathsf{rk}^{\pi,\mathsf{next}}}}$ and the parties invoke \cref{piop:rotcheck} on $(\mathsf{rk}^{\pi}, \mathsf{rk}^{\pi,\mathsf{next}}, -1)$, then the $\mathsf{NonNeg}$ \cref{piop:sign-check} on the column with activator $\contig_{\mu,n-1}$ and data $\mathsf{rk}^{\pi,\mathsf{next}} - \mathsf{rk}^{\pi}$.
    \item (\emph{Strings sorted})\label[step]{rc:lex} $\PIOPProver$ and $\PIOPVerifier$ invoke \cref{piop:lex-sorted-check} on $\strcol^{\pi}$, obtaining its equal-neighbour indicator $\mathsf{same}^{\pi}$, which marks the adjacent pairs with $s^{\pi}_{i} = s^{\pi}_{i+1}$.
    \item (\emph{Ties coincide})\label[step]{rc:ties} $\PIOPProver$ and $\PIOPVerifier$ invoke \cref{piop:zerocheck} on $\mathsf{same}^{\pi} \mult (\mathsf{rk}^{\pi,\mathsf{next}} - \mathsf{rk}^{\pi})$, and the $\mathsf{Pos}$ \cref{piop:sign-check} on the column with activator $\contig_{\mu,n-1} \mult (1 - \mathsf{same}^{\pi})$ and data $\mathsf{rk}^{\pi,\mathsf{next}} - \mathsf{rk}^{\pi}$.
  \end{enumerate}
\end{DefinePIOP}

\Cref{rc:mono} is implied by \cref{rc:ties} and is listed only for clarity; it may be dropped. Density of the ranks (that they are exactly $0, \dots, D-1$ for $D$ distinct strings) is not needed for any protocol and is not checked; an honest data owner produces dense ranks, and any rank column in $\NPRelation_{\mathsf{RankValid}}$ orders the strings correctly.

\parhead{Sorting by rank}
With a trusted rank column in hand, the \texttt{ORDER BY} operator is proven exactly as in \cref{piop:str-sort}, except that the char-level \cref{ss:order} is replaced by an integer sortedness check on the inherited ranks: by $\NPRelation_{\mathsf{RankValid}}$, non-decreasing ranks along the output are equivalent to $s_i \lexle s_{i+1}$.

\begin{DefinePIOP}{Rank-Based String Sort Check}{rank-sort}[Rank-Based String Sort Check]
  \emergencystretch=3em\relax
  \emph{Inputs:} the input column $\Oracle{\MLE{\strcol^{\mathsf{old}}}}$ and the output column $\Oracle{\MLE{\strcol^{\mathsf{new}}}}$, both carrying rank oracles $\Oracle{\MLE{\mathsf{rk}^{\mathsf{old}}}}, \Oracle{\MLE{\mathsf{rk}^{\mathsf{new}}}}$; and $\Oracle{\MLE{\mathsf{ind}}}$.
  \vspace{2mm}
  \begin{enumerate}[nosep]
    \item (\emph{Validity, provenance, permutation})\label[step]{rs:valid} As in \cref{ss:valid,ss:perm}; the provenance lookup of \cref{piop:domain-preserving-update-check} transfers $\mathsf{rk}^{\mathsf{old}}$ to $\mathsf{rk}^{\mathsf{new}}$.
    \item (\emph{Order})\label[step]{rs:order} $\PIOPProver$ sends $\Oracle{\MLE{\mathsf{rk}^{\mathsf{new},\mathsf{next}}}}$ and the parties invoke \cref{piop:rotcheck} on $(\mathsf{rk}^{\mathsf{new}}, \mathsf{rk}^{\mathsf{new},\mathsf{next}}, -1)$, then the $\mathsf{NonNeg}$ \cref{piop:sign-check} on the column with activator $\contig_{\mu,n-1}$ and data $\mathsf{rk}^{\mathsf{new},\mathsf{next}} - \mathsf{rk}^{\mathsf{new}}$.
  \end{enumerate}
\end{DefinePIOP}

The online cost of the string sort is thus that of $\truthtable$'s integer sort, independent of the string lengths and of the char-level columns, which the planner may already have dropped (\cref{sec:string-arith}). Multi-key sorts and descending order carry over verbatim from \cref{sec:lex-sorted-check}, with the tie indicator $\mathbf{1}[\mathsf{rk}^{\mathsf{new},\mathsf{next}} = \mathsf{rk}^{\mathsf{new}}]$, since equal ranks are equal strings.

\parhead{Scope}
Ranks exist only for strings that descend from a committed column. A sort key produced by a \emph{domain-altering} expression (\texttt{SUBSTRING}, \texttt{CONCAT}, case mapping) is a new string with no rank and falls back to \cref{piop:lex-sorted-check}; the same holds for comparisons against a literal that need not occur in the column. In TPC-H, nine of the ten queries that order by a string column sort on copies of base columns (Q1, Q2, Q4, Q7, Q9, Q12, Q16, Q20, Q21) and are covered by \cref{piop:rank-sort}. The exception is Q22, whose key $\texttt{cntrycode} = \texttt{SUBSTRING}(\texttt{c\_phone}, 1, 2)$ is a freshly constructed string: \texttt{SUBSTRING} is a domain-altering expression outside the scope of our string protocols, so neither the rank path nor \cref{piop:lex-sorted-check} applies to it directly, as there is no trusted arithmetization of the new column to apply them to. A fixed-length prefix, however, needs no string machinery at all: the two projections $\mathsf{c}_0, \mathsf{c}_1$ of \texttt{c\_phone} (\cref{piop:char-proj}) pack into the single integer $\mathsf{c}_0 \mult 2^{\alpha} + \mathsf{c}_1$, an injective and order-preserving encoding of the key, so the \texttt{GROUP BY}, the \texttt{IN} list and the \texttt{ORDER BY} of Q22 all become integer operations of $\truthtable$, and the verifier decodes the integer back into two characters when opening the result.\ali{Verify the query list against the TPC-H specification, and whether $\truthtable$ dropped these \texttt{ORDER BY} clauses.}
